\chapter{AR, MA e ARMA}
\label{cap:arma}
% ── índice ──────────────────────────────────────────────────────
\index{ARMA|textbf}
\index{autoregressivo|textbf}
\index{média móvel|textbf}
\index{arima@\texttt{arima()}|textbf}
\index{autocorrelação|textbf}
\index{AIC|textbf}
\index{BIC|textbf}
\index{Box-Jenkins, metodologia de|textbf}
\index{estacionariedade|textbf}

% ─────────────────────────────────────────────────────────────────────────────
\section{Processos estocásticos e estacionariedade}

Um \textbf{processo estocástico} $\{y_t\}$ é estacionário (em covariância)
quando:
\[
  E[y_t] = \mu, \qquad
  \mathrm{Var}[y_t] = \sigma^2, \qquad
  \mathrm{Cov}(y_t, y_{t-k}) = \gamma_k
\]
para todo $t$ — média, variância e autocovariâncias dependem apenas da
defasagem $k$, não do tempo $t$.

O caso mais simples é o \textbf{ruído branco} $\varepsilon_t \sim
\mathrm{WN}(0, \sigma^2)$: $E[\varepsilon_t] = 0$, $\mathrm{Var}[\varepsilon_t]
= \sigma^2$ e $\mathrm{Cov}(\varepsilon_t, \varepsilon_s) = 0$ para $t \neq
s$. Ruído branco é o bloco de construção de todos os modelos desta seção.

% ─────────────────────────────────────────────────────────────────────────────
\section{Processo AR(\textit{p})}

O modelo autorregressivo de ordem $p$ (\textbf{AR($p$)}) expressa $y_t$
como combinação linear de seus próprios valores passados mais um choque:

\[
  y_t = c + \phi_1 y_{t-1} + \phi_2 y_{t-2} + \cdots + \phi_p y_{t-p}
        + \varepsilon_t
\]

Para o AR(1), a condição de estacionariedade é $|\phi_1| < 1$. Quando
satisfeita, a média incondicional é $\mu = c / (1 - \phi_1)$ e a
autocorrelação decai geometricamente:
\[
  \rho_k = \phi_1^k, \qquad k = 0, 1, 2, \ldots
\]

A \textbf{Função de Autocorrelação Parcial} (PACF) corta em zero após
defasagem $p$ — característica diagnóstica do AR.

% ─────────────────────────────────────────────────────────────────────────────
\section{Processo MA(\textit{q})}

O modelo de médias móveis de ordem $q$ (\textbf{MA($q$)}) expressa $y_t$
como combinação linear de choques presentes e passados:

\[
  y_t = \mu + \varepsilon_t + \theta_1 \varepsilon_{t-1}
        + \cdots + \theta_q \varepsilon_{t-q}
\]

O MA($q$) é sempre estacionário. A \textbf{Função de Autocorrelação} (ACF)
corta em zero após defasagem $q$ — em contraste com o AR, cuja ACF decai
geometricamente. A condição de \textbf{invertibilidade} ($|\theta_1| < 1$
para MA(1)) garante representação AR($\infty$) única.

\begin{nota}
  Diagnóstico padrão: ACF decai lentamente e PACF corta → AR($p$);
  ACF corta e PACF decai lentamente → MA($q$); ambas decaem → ARMA.
\end{nota}

% ─────────────────────────────────────────────────────────────────────────────
\section{Processo ARMA(\textit{p},\,\textit{q})}

O ARMA($p$,$q$) combina os dois mecanismos:
\[
  y_t = c + \sum_{i=1}^p \phi_i y_{t-i}
        + \varepsilon_t + \sum_{j=1}^q \theta_j \varepsilon_{t-j}
\]

É estacionário quando as raízes do polinômio AR estão fora do círculo
unitário e invertível quando as raízes do polinômio MA também estão.

% ─────────────────────────────────────────────────────────────────────────────
\section{Verificação por DGP}

Como não existe gerador de defasagem em \lstinline{generate}, os processos
são construídos com um laço \lstinline{while} que preenche $y_t$ linha a
linha, usando \lstinline{mean(..., if = t == k)} como acessor de linha $k$.

\subsection*{AR(1) com $\phi = 0{,}7$}

\begin{lstlisting}[caption={DGP AR(1) e estimação}]
seed(42)
input df
id
1
end
for i in 1..=9 { df = append(df, df) }
df |> filter(_n <= 500)
generate df e = rnormal()
generate df y = 0.0
generate df t = _n
let n = nrow(df)
let i = 2
while i <= n {
    let lag_y = mean(df, y, if = t == i - 1)
    let et    = mean(df, e, if = t == i)
    replace df y = 0.7 * lag_y + et if t == i
    i = i + 1
}
let m = arima(df, y, p=1, d=0, q=0)
display m
\end{lstlisting}

\begin{lstlisting}[language={}, basicstyle=\ttfamily\small,
  frame=none, numbers=none, backgroundcolor=\color{codbg},
  xleftmargin=0pt, xrightmargin=0pt, breaklines=false]
================== ARIMA(1,0,0) via Hannan-Rissanen ==================
Observations:               494
AIC:                   157.7072   BIC:                   166.1123
Parameters:
  intercept   0.5454   std=0.0557   z=9.786   p=0.000   [0.4362, 0.6547]
  ar.L1       0.6638   std=0.0335   z=19.814  p=0.000   [0.5981, 0.7295]
======================================================================
\end{lstlisting}

$\hat\phi_1 = 0{,}6638$ com verdadeiro $\phi_1 = 0{,}7$; o IC 95\% $[0{,}60;\;
0{,}73]$ cobre o valor verdadeiro.

\subsection*{MA(1) com $\theta = 0{,}5$}

\begin{lstlisting}[caption={DGP MA(1) e estimação}]
generate df y = e
let i = 2
while i <= n {
    let lag_e = mean(df, e, if = t == i - 1)
    replace df y = y + 0.5 * lag_e if t == i
    i = i + 1
}
let m = arima(df, y, p=0, d=0, q=1)
display m
\end{lstlisting}

\begin{lstlisting}[language={}, basicstyle=\ttfamily\small,
  frame=none, numbers=none, backgroundcolor=\color{codbg},
  xleftmargin=0pt, xrightmargin=0pt, breaklines=false]
================== ARIMA(0,0,1) via Hannan-Rissanen ==================
Observations:               493
AIC:                   162.5430   BIC:                   170.9441
Parameters:
  intercept   0.7297   std=0.0128   z=57.014  p=0.000   [0.7046, 0.7548]
  ma.L1       0.5087   std=0.0456   z=11.163  p=0.000   [0.4194, 0.5980]
======================================================================
\end{lstlisting}

$\hat\theta_1 = 0{,}5087 \approx 0{,}5$; IC 95\% $[0{,}42;\; 0{,}60]$
cobre $\theta_1 = 0{,}5$.

\subsection*{ARMA(1,1) com $\phi=0{,}6$ e $\theta=0{,}4$}

\begin{lstlisting}[caption={DGP ARMA(1{,}1) e estimação}]
generate df y = e
let i = 2
while i <= n {
    let lag_y = mean(df, y, if = t == i - 1)
    let lag_e = mean(df, e, if = t == i - 1)
    let et    = mean(df, e, if = t == i)
    replace df y = 0.6*lag_y + et + 0.4*lag_e if t == i
    i = i + 1
}
let m = arima(df, y, p=1, d=0, q=1)
display m
\end{lstlisting}

\begin{lstlisting}[language={}, basicstyle=\ttfamily\small,
  frame=none, numbers=none, backgroundcolor=\color{codbg},
  xleftmargin=0pt, xrightmargin=0pt, breaklines=false]
================== ARIMA(1,0,1) via Hannan-Rissanen ==================
Observations:               493
AIC:                   159.4217   BIC:                   172.0233
Parameters:
  intercept   0.8073   std=0.0660   z=12.227  p=0.000   [0.6779, 0.9367]
  ar.L1       0.5259   std=0.0380   z=13.831  p=0.000   [0.4514, 0.6004]
  ma.L1       0.4829   std=0.0591   z=8.167   p=0.000   [0.3670, 0.5987]
======================================================================
\end{lstlisting}

Ambos os parâmetros são recuperados com boa precisão: $\hat\phi = 0{,}53$
(verdadeiro $0{,}6$) e $\hat\theta = 0{,}48$ (verdadeiro $0{,}4$).

% ─────────────────────────────────────────────────────────────────────────────
\section{Seleção de ordem}

O \hayashi{} reporta AIC e BIC automaticamente. O critério prático:

\begin{enumerate}
  \item Estime modelos com ordens $(p,q)$ em uma grade pequena (ex.\ $0
        \leq p,q \leq 3$).
  \item Selecione o modelo com menor AIC (preferência por ajuste) ou BIC
        (preferência por parcimônia).
  \item Verifique que os resíduos são ruído branco.
\end{enumerate}

\begin{lstlisting}[caption={Seleção de ordem por AIC}]
let m00 = arima(df, y, p=0, d=0, q=0)
let m10 = arima(df, y, p=1, d=0, q=0)
let m01 = arima(df, y, p=0, d=0, q=1)
let m11 = arima(df, y, p=1, d=0, q=1)
esttab(m00, m10, m01, m11)
\end{lstlisting}

% ─────────────────────────────────────────────────────────────────────────────
\section{Sintaxe}

\begin{lstlisting}
arima(df, y, p=1, d=0, q=0)
arima(df, y, p=0, d=0, q=1)
arima(df, y, p=1, d=1, q=1)
\end{lstlisting}

\begin{center}
\begin{tabular}{lll}
\toprule
\textbf{Parâmetro} & \textbf{Padrão} & \textbf{Descrição} \\
\midrule
\texttt{p} & \texttt{0} & ordem autorregressiva \\
\texttt{d} & \texttt{0} & grau de diferenciação \\
\texttt{q} & \texttt{0} & ordem de médias móveis \\
\bottomrule
\end{tabular}
\end{center}

\begin{lstlisting}[caption={Fluxo típico com dados reais}]
load "ipca.csv" as df

let m = arima(df, inflacao, p=1, d=0, q=1)
display m
let yhat = predict(m, "xb")
\end{lstlisting}

\begin{nota}
  O laço \lstinline{while} dos DGPs acima é $O(n^2)$ e serve apenas para
  ilustrar a estrutura do processo. Em uso real, os dados são carregados
  de arquivo com \lstinline{load} e o laço não é necessário.
\end{nota}
